\subsection{强凸性与光滑性}\label{sec:A.2}

\begin{definition}\label{def:A.7}
若函数$f:\mathbb R^d\to\mathbb R\cup\{\infty\}$对任意$x,y\in\mathbb R^d$及$t\in[0,1]$满足
$$
f((1-t)x+ty)\leq(1-t)f(x)+tf(y)-\frac{\alpha t(1-t)}2\|y-x\|^2,
$$
则称$f$为$\alpha$-凸的（$\alpha$-convex），其中$\alpha\in\mathbb R$.
\end{definition}
$\alpha=0$的情形对应于\autoref{def:A.1}的凸性. 当$\alpha>0$时，称$f$为强凸的（strongly convex），上式加强了通常的凸性不等式；当$\alpha<0$时，称$f$为半凸的（semi-convex）.

当$f$连续可微时，$\alpha$-凸性等价于下面第一条陈述. 当$f$二阶连续可微时，$\alpha$-凸性与下面两条陈述都等价.
\begin{enumerate}
\item 对任意$x,y\in\mathbb R^d$，$f(y)\geq f(x)+\langle\nabla f(x),y-x\rangle+\frac\alpha2\|x-y\|^2$.
\item 对任意$x\in\mathbb R^d$，$\nabla^2f(x)\succeq\alpha I$.
\end{enumerate}
我们也给出与之对偶的性质，即二阶导数上界.

\begin{definition}\label{def:A.8}
若连续可微函数$f:\mathbb R^d\to\mathbb R$对任意$x,y\in\mathbb R^d$满足
$$
f(y)\leq f(x)+\langle\nabla f(x),y-x\rangle+\frac\beta2\|y-x\|^2,
$$
则称$f$为$\beta$-光滑的（$\beta$-smooth），其中$\beta\geq0$.
\end{definition}

当$f$二阶连续可微时，这等价于对任意$x\in\mathbb R^d$有$\nabla^2f(x)\preceq\beta I$. 若$f$还为凸函数，则$\nabla^2f(x)\succeq0$，因此$\nabla^2f(x)$的算子范数至多为$\beta$. 这又等价于映射$\nabla f:\mathbb R^d\to\mathbb R^d$是$\beta$-Lipschitz的（因此，凸且光滑的函数常被称为梯度Lipschitz函数）.

强凸性与光滑性互为对偶性质. 为简单起见，下面假设$f$连续可微；这些假设其实可以适当放宽.

\begin{lemma}\label{lem:A.9}
设$f:\mathbb R^d\to\mathbb R$连续可微、凸，并且当$\|x\|\to\infty$时有$\|\nabla f(x)\|\to\infty$. 设$\alpha>0$. 则$f$为$\alpha$-强凸函数，当且仅当其凸共轭$f^*$为$1/\alpha$-光滑函数.
\end{lemma}

\begin{proof}
由凸分析中的经典结果（参见\cite[定理25.5、定理26.6和引理26.7]{ref274}），在我们的假设下，$\nabla f:\mathbb R^d\to\mathbb R^d$是微分同胚，其逆为$\nabla f^*$.

（$\implies$）将强凸性的一阶条件与交换$x,y$所得的不等式相加，得到
$$
\langle\nabla f(x)-\nabla f(y),x-y\rangle\geq\alpha\|x-y\|^2.
$$
令$x=\nabla f^*(x')$、$y=\nabla f^*(y')$，并回顾$\nabla f\circ\nabla f^*=\mathrm{id}$. 上式说明，对任意$x',y'\in\mathbb R^d$，
$$
\langle x'-y',\nabla f^*(x')-\nabla f^*(y')\rangle
\geq\alpha\|\nabla f^*(x')-\nabla f^*(y')\|^2.
$$
对左侧应用Cauchy--Schwarz不等式并整理，可知$\nabla f^*$是$1/\alpha$-Lipschitz的. 如上所述，这等价于$f^*$为$1/\alpha$-光滑函数.

（$\impliedby$）由$f^*$的光滑性，
$$
\begin{aligned}
f(y)&=\sup_{y'\in\mathbb R^d}\{\langle y,y'\rangle-f^*(y')\}\\
&\geq\sup_{y'\in\mathbb R^d}\bigl\{\langle y,y'\rangle-f^*(x')
-\langle\nabla f^*(x'),y'-x'\rangle-\tfrac1{2\alpha}\|y'-x'\|^2\bigr\}\\
&=-f^*(x')+\langle y,x'\rangle+\frac\alpha2\|y-\nabla f^*(x')\|^2.
\end{aligned}
$$
取$x'=\nabla f(x')$，使得$\nabla f^*(x')=x$，并回顾$f(x)+f^*(x')=\langle x,x'\rangle$. 于是
$$
f(y)-f(x)-\langle\nabla f(x),y-x\rangle\geq\frac\alpha2\|y-x\|^2,
$$
证明完成.
\end{proof}

注意，若$f,f^*$二阶连续可微，则$f$为$\alpha$-强凸函数当且仅当$\nabla^2 f\succeq\alpha I$；而$f^*$为$1/\alpha$-光滑函数当且仅当
$$
\nabla^2f^*=(\nabla^2f)^{-1}\circ\nabla f^*\preceq\alpha^{-1}I,
$$
这给出了一个更直观的证明.

强凸性还蕴含下面的性质；它可以视为“凸性保证局部最优点为全局最优点”这一原理的一个强定量形式.

\begin{definition}[Polyak--Łojasiewicz不等式]\label{def:A.10}\index{Polyak–Łojasiewicz (PŁ) inequality@Polyak–Łojasiewicz（PŁ）不等式（Polyak–Łojasiewicz (PŁ) inequality）|hyperpage}
若连续可微函数$f:\mathbb R^d\to\mathbb R$对任意$x\in\mathbb R^d$满足
$$
\|\nabla f(x)\|^2\geq2\alpha\bigl(f(x)-\inf f\bigr),
$$
则称$f$满足常数为$\alpha>0$的Polyak--Łojasiewicz（PŁ）不等式（Polyak--Łojasiewicz inequality）.
\end{definition}

\begin{lemma}\label{lem:A.11}
若$f:\mathbb R^d\to\mathbb R$连续可微且为$\alpha$-凸函数，其中$\alpha>0$，则它满足常数为$\alpha$的PŁ不等式.
\end{lemma}
\begin{proof}
以$x_\star$表示$f$的极小化点. 则
$$
\inf f=f(x_\star)\geq f(x)+\langle\nabla f(x),x_\star-x\rangle
+\frac\alpha2\|x_\star-x\|^2.
$$
由Cauchy--Schwarz不等式和Young不等式，
$$
\begin{aligned}
\langle\nabla f(x),x_\star-x\rangle
&\geq-\|\nabla f(x)\|\,\|x_\star-x\|\\
&\geq-\frac1{2\alpha}\|\nabla f(x)\|^2-\frac\alpha2\|x_\star-x\|^2.
\end{aligned}
$$
代入并整理即得证明.
\end{proof}

还要指出，强凸性蕴含函数在极小化点附近的二次增长.
\begin{lemma}\label{lem:A.12}
若$f:\mathbb R^d\to\mathbb R\cup\{\infty\}$为$\alpha$-凸函数，$x_\star$为$f$的极小化点，则对任意$x\in\mathbb R^d$，
\begin{equation}\label{eq:A.2}
f(x)-f(x_\star)\geq\frac\alpha2\|x-x_\star\|^2.
\end{equation}
\end{lemma}
\begin{proof}
强凸性不等式给出
$$
f((1-t)x_\star+tx)\leq(1-t)f(x_\star)+tf(x)
-\frac{\alpha t(1-t)}2\|x-x_\star\|^2,
$$
或者写成
$$
\begin{aligned}
0&\leq f((1-t)x_\star+tx)-f(x_\star)\\
&\leq t\left[f(x)-f(x_\star)-\frac{\alpha(1-t)}2\|x-x_\star\|^2\right].
\end{aligned}
$$
整理，两边除以$t$，再令$t\searrow0$，即证结论.
\end{proof}

事实上，在\autoref{sec:5}的\autoref{exer:5.5}中，我们将进一步说明，PŁ不等式本身就蕴含增长不等式\eqref{eq:A.2}. 关于这些不等式之间关系的简明论述，参见\cite[附录A]{ref174}.
